\chapter{第 16 章分层答案}
\label{app:ch16-solutions}

\section*{口述概念}
\begin{enumerate}[leftmargin=*]
  \item $t_e$ 是事件发生，$t_a$ 是值首次可用，$t_d$ 是形成决策，$t_y$ 是标签开始。$t_a\le t_d$ 阻止使用尚未可得的数据，$t_d<t_y$ 阻止特征窗口进入未来标签；完整协议还要求 $t_e\le t_a$。
  \item 测试集的职责是评估冻结流程。研究者若根据测试结果修改特征、窗口、阈值或叙事，测试信息就参与了选择；后续得分是开发证据，不再是一次性确认性证据。应另留最终区间，或诚实地重新定义研究阶段。
  \item 复现问“同一输入能否得到同一输出”；统计有效问“在声明假设下推断是否控制误差”；可交易问“订单在约束、流动性和成本下能否实现”。可复现的泄漏回测仍无效，统计显著的纸面信号也可能无法成交。
\end{enumerate}

\section*{笔试推导}
\begin{enumerate}[leftmargin=*]
  \item 独立真零假设下，每次不误拒概率为 $1-\alpha$，全部不误拒为 $(1-\alpha)^m$，所以至少一次误拒为 $1-(1-\alpha)^m$。代入得 $1-0.95^{20}\approx0.6415$。
  \item 令 $A_j=\{p_j\le\alpha/m\}$。并集界给出
  $\Pr(\cup_jA_j)\le\sum_j\Pr(A_j)\le m(\alpha/m)=\alpha$，不需要检验独立。阈值为 $0.0025$，所以 $0.002$ 通过、$0.003$ 拒绝。
  \item $R^{\mathrm{gross}}=0.04$。单笔成本为 $0.002$，三笔总成本 $C=0.006$，故 $R^{\mathrm{net}}=0.034$。若观察是连续持有期收益，应复利；若资本权重不同，应先加权；若交易重叠，还需处理同时占用资本与净额成交。
\end{enumerate}

\section*{数值编程}
\begin{enumerate}[leftmargin=*]
  \item 解析 ISO 时间后逐行检查两个复合条件。正常行必须满足 $t_e\le t_a\le t_d$ 且 $t_d<t_y$；把第一行 $t_a$ 改为 9:31，而 $t_d$ 保持 9:25，应稳定失败为“row 1 was not available before the decision”。不要把当前日期或机器时区作为 oracle。
  \item 把每个样本标签表示为半开区间 $[s_i,e_i)$。验证折为 $V$ 时，从训练候选中删除所有满足
  \[
    [s_i,e_i)\cap[s_j,e_j)\ne\varnothing,\qquad j\in V
  \]
  的样本。反例可取相邻两日信号都使用未来五日收益：随机分到两折后仍共享四个未来价格。
  \item 为每个订单记录决策基准价、可成交数量、实际成交均价和期末基准价。成交部分扣佣金、价差与规模冲击；未成交部分的机会成本按预先声明的基准计算。分别报告 executed-only P\&L 与 implementation shortfall，不能把未成交订单当作零成本成交。
\end{enumerate}

\section*{研究判断}
\begin{enumerate}[leftmargin=*]
  \item 缺失证据包括：完整 500 次及其变体的尝试账本；预先声明的错误率控制；选择过程被纳入重抽样或嵌套切分；完全未参与选择的最终时间外测试；成本、容量和泄漏审计。可把探索期与确认期分开，在开发折内完成筛选和超参数选择，只在冻结后打开最终测试，并同时报告赢家与搜索规模。
  \item 可能改变成交结论的包括：买入后可售时间、涨跌幅和有效申报价格、停复牌、交易日历、复权可得时间、融资融券标的与券源。条件期望、误差控制、样本外定义和成本账本等数学原则仍通用，但输入约束与可行集合必须按市场和研究期重写。
  \item 应拒绝“十倍资金仍保持原净表现”和据此得到的容量结论。需要按订单规模、成交量参与率、波动率与执行时长估计冲击，报告盘口/成交量覆盖、未成交、实现差额、压力情景和净绩效随资金规模的曲线，并说明模型不能覆盖的极端流动性状态。
\end{enumerate}

\section*{代码恢复路径}
在仓库根目录运行
\begin{center}
\path{uv run jupyter nbconvert --to notebook --execute --ExecutePreprocessor.allow_error_names=SystemExit notebooks/upper/ch16_research_validity.ipynb}
\end{center}
正常路径必须同时给出 timeline、timezone、calendar、split、multiplicity、performance 和 friction 七类通过证据。随后复制 oracle：先把第一行 \path{available_at} 改到 \path{decision_at} 之后，确认时间门禁失败；再恢复时间并把 \path{selected_raw_p_value} 改为 $0.003$，确认多重检验门禁失败。把任一时间改成错误 UTC offset 或把决策日移出 \path{2024-07-v1} 日历，也必须稳定失败。恢复完成的标准不是“改回能跑”，而是能解释负例分别破坏了哪个研究假设，以及为什么增加样本或更换随机 seed 都无法修复。

\section*{独立 oracle 证据对照}
本章的发布证据固定为：\textbf{事件时间、可得时间、决策时间与收益起点}满足先后约束；Bonferroni 阈值为 $0.0025$；三笔交易的\textbf{总成本为 $0.006$}；\textbf{成本后收益为 $0.034$}。市场条件协议要求逐研究期核验\textbf{T+1、涨跌停成交、停牌、复权可得性与融券资格}，而不是把当前规则参数当作跨时期常数。
\section*{逐题逐步核对}
\begin{enumerate}[leftmargin=*]
  \item 事件时间 $t_e$、首次可得 $t_a$、决策时间 $t_d$、标签起点 $t_y$ 必须满足 $t_e\le t_a\le t_d<t_y$。逐行解析时区和日历，任何一条违反都拒绝；哈希只能证明字节一致，不能证明时间语义正确。
  \item 测试集只能评估冻结流程；若根据测试分数改特征、阈值、窗口或叙事，测试已进入选择集。保留未见最终区间，或把结果明确降级为开发/探索证据；不能继续称一次性样本外。
  \item 复现、统计有效性和可交易性分别检查字节/命令、误差控制、订单/成本约束。泄漏回测可以完全可复现，显著结果也可能无法成交；答案要把三种结论分栏。
  \item 十次独立零假设查看的 FWER 为 $1-(1-.05)^{10}=.401263$；Bonferroni 用 $\alpha/20=.0025$，所以 $p=.002$ 通过而 .003 不通过。先把检验族、alpha-spending 和停止规则写进协议，再读取数据。
  \item 训练/验证/测试区间用标签半开区间 $[s_i,e_i)$ 表示；验证样本 $j$ 存在时，删除所有与 $[s_j,e_j)$ 相交的训练样本，并在验证末端加 embargo。逐行打印保留索引，才能确认 purge 不是只改了标题。
  \item 三笔交易总毛收益 .04；每笔成本 .002，$C=.006$，净收益 .034。未成交部分不能当作零成本成交，应分 executed-only PnL 和 implementation shortfall，并记录 arrival price、成交量、价差和冲击。
  \item 多重检验审计必须保留完整候选、p 值、选择规则、校正方法和未通过项；A 股还需逐研究期检查 T+1、涨跌停、停牌、复权可得性、融券资格和成交完成率。当前市场规则不能自动外推到历史期。
  \item 容量报告要把资金规模、参与率、波动、执行时长和冲击模型放入函数，画净绩效/容量曲线并报告未成交和压力场景；“资金增十倍仍相同表现”没有订单级证据时应拒绝。
\end{enumerate}
